\section{Conclusion}
\label{sec:conclusion}

We presented the \textbf{Specification-guided Defect Prevention Framework}
(\textsc{SgDP}), a general architecture for preventing specification-conformance defects
in LLM-generated implementations.  The framework grounds every component in specification
semantics: witnesses are generated by encoding guard precedence as SMT constraints;
repair feedback is structured by a typed Semantic Feedback IR carrying specification-level
information; acceptance is gated by a conjunctive predicate combining oracle-grounded
conformance with specification-derived quality criteria.

HSP-Agile reaches 86.3\% mean strict formal conformance at seed~0 on canonical
E1; B1 leads at 89.2\%.
Mechanism analysis confirms that gains are overlap-conditional, not universal:
E3 shows the advantage grows with specification complexity; E4 shows SMT witnesses
maintain boundary coverage where random sampling fails; E5/E6 show that the Semantic
Feedback IR produces faster repair convergence than test-only or test-plus-expected
feedback.  Three formal theorems establish the soundness, safety, and monotonicity
properties that underpin these empirical results.

Generalisation experiments (E8) demonstrate that the framework advantage holds across
Mini-StateMachine and Mini-Z notations with only adapter-level changes, supporting the
framework claim rather than an Agile-SOFL-specific claim.  Failure analysis (E9)
identifies the remaining failure modes — primarily ordering-repair failures and
arithmetic-boundary cases — and guides the most promising future extensions.

Future work includes extending the Z3 encoding to non-linear arithmetic guards,
scaling the pattern catalogue via automated mining of specification fault databases,
evaluating on industrial Agile-SOFL projects, and exploring adaptive repair strategies
that prioritise counterexamples by fault severity based on the IR's priority field.
